% Figure: The Attention Gap — Google Trends for agent terminology (2004-2026)
% For inclusion in Section 1 (Introduction) of sutra-main paper
% Data source: Google Trends, quarterly averages
%
% Usage: % Figure: The Attention Gap — Google Trends for agent terminology (2004-2026)
% For inclusion in Section 1 (Introduction) of sutra-main paper
% Data source: Google Trends, quarterly averages
%
% Usage: % Figure: The Attention Gap — Google Trends for agent terminology (2004-2026)
% For inclusion in Section 1 (Introduction) of sutra-main paper
% Data source: Google Trends, quarterly averages
%
% Usage: % Figure: The Attention Gap — Google Trends for agent terminology (2004-2026)
% For inclusion in Section 1 (Introduction) of sutra-main paper
% Data source: Google Trends, quarterly averages
%
% Usage: \input{figures/search-trends.tex}

\begin{figure*}[t]
\centering
\begin{tikzpicture}
\begin{axis}[
    width=\textwidth,
    height=7.5cm,
    xlabel={Year},
    ylabel={Google Trends Interest (0--100)},
    xmin=2004, xmax=2026.25,
    ymin=-3, ymax=108,
    xtick={2004,2006,2008,2010,2012,2014,2016,2018,2020,2022,2024,2026},
    xticklabel style={font=\small},
    yticklabel style={font=\small},
    xlabel style={font=\small},
    ylabel style={font=\small},
    legend style={
        at={(0.02,0.98)},
        anchor=north west,
        font=\small,
        draw=black!30,
        fill=white,
        fill opacity=0.9,
        text opacity=1,
        cells={anchor=west},
    },
    grid=major,
    grid style={gray!20},
    clip=false,
    % Era shading
    extra x ticks={},
]

% ── The Dead Zone shading (2013-2022) ──────────────────────
\fill[black!5] (axis cs:2012.5,-3) rectangle (axis cs:2022.5,108);
\node[font=\footnotesize\itshape, text=black!40, rotate=90]
    at (axis cs:2017.5,50) {The Attention Gap};

% ── Key event markers ─────────────────────────────────────
% AlexNet
\draw[black!40, dashed, thin] (axis cs:2012,0) -- (axis cs:2012,105);
\node[font=\tiny, text=black!50, anchor=south, rotate=45]
    at (axis cs:2012,105) {AlexNet};

% Transformers
\draw[black!40, dashed, thin] (axis cs:2017,0) -- (axis cs:2017,105);
\node[font=\tiny, text=black!50, anchor=south, rotate=45]
    at (axis cs:2017,105) {Transformers};

% GPT-3
\draw[black!40, dashed, thin] (axis cs:2020,0) -- (axis cs:2020,105);
\node[font=\tiny, text=black!50, anchor=south, rotate=45]
    at (axis cs:2020,105) {GPT-3};

% GPT-4 / AutoGen
\draw[black!40, dashed, thin] (axis cs:2023.25,0) -- (axis cs:2023.25,105);
\node[font=\tiny, text=black!50, anchor=south, rotate=45]
    at (axis cs:2023.25,105) {GPT-4};

% ── Data series ────────────────────────────────────────────
% "Intelligent agent" — the dominant new term (area fill + line)
\addplot[
    color=blue!70!black,
    fill=blue!10,
    thick,
    smooth,
    mark=none,
] table[x index=0, y index=3] {figures/search-trends.dat}
    \closedcycle;

% "Agentic AI" — the new buzzword
\addplot[
    color=red!70!black,
    fill=red!10,
    thick,
    smooth,
    mark=none,
] table[x index=0, y index=2] {figures/search-trends.dat}
    \closedcycle;

% "Multi-agent system" — the classical term
\addplot[
    color=black!70!green,
    fill=green!8,
    thick,
    smooth,
    mark=none,
] table[x index=0, y index=1] {figures/search-trends.dat}
    \closedcycle;

\legend{
    ``Intelligent agent'',
    ``Agentic AI'',
    ``Multi-agent system''
}

% ── Peak annotation ────────────────────────────────────────
\node[font=\tiny, anchor=south west, text=blue!70!black]
    at (axis cs:2025.3,88) {100 (Oct 2025)};

% ── Classical era label ────────────────────────────────────
\draw[decorate, decoration={brace, amplitude=4pt, mirror},
      black!50, thick]
    (axis cs:2004,-2) -- (axis cs:2006,-2)
    node[midway, below=5pt, font=\tiny, text=black!50]
    {Classical MAS};

\end{axis}
\end{tikzpicture}
\caption{\textbf{The Attention Gap.} Google Trends search interest
(2004--2026) for three agent-related terms. ``Multi-agent system''
(the classical terminology) peaked at 3--4 in 2004 and flatlined near
zero by 2012. ``Intelligent agent'' and ``Agentic AI'' (new terminology)
exploded starting Q4 2024, reaching 100 by October 2025. The shaded
region marks the Attention Gap (2013--2022): a decade where agent research
effectively disappeared from public discourse. The new wave of interest
uses different terminology than the classical research it unknowingly
reinvents---the terminological disconnect that motivates this survey.
Data source: Google Trends, quarterly averages.}
\label{fig:trends}
\end{figure*}


\begin{figure*}[t]
\centering
\begin{tikzpicture}
\begin{axis}[
    width=\textwidth,
    height=7.5cm,
    xlabel={Year},
    ylabel={Google Trends Interest (0--100)},
    xmin=2004, xmax=2026.25,
    ymin=-3, ymax=108,
    xtick={2004,2006,2008,2010,2012,2014,2016,2018,2020,2022,2024,2026},
    xticklabel style={font=\small},
    yticklabel style={font=\small},
    xlabel style={font=\small},
    ylabel style={font=\small},
    legend style={
        at={(0.02,0.98)},
        anchor=north west,
        font=\small,
        draw=black!30,
        fill=white,
        fill opacity=0.9,
        text opacity=1,
        cells={anchor=west},
    },
    grid=major,
    grid style={gray!20},
    clip=false,
    % Era shading
    extra x ticks={},
]

% ── The Dead Zone shading (2013-2022) ──────────────────────
\fill[black!5] (axis cs:2012.5,-3) rectangle (axis cs:2022.5,108);
\node[font=\footnotesize\itshape, text=black!40, rotate=90]
    at (axis cs:2017.5,50) {The Attention Gap};

% ── Key event markers ─────────────────────────────────────
% AlexNet
\draw[black!40, dashed, thin] (axis cs:2012,0) -- (axis cs:2012,105);
\node[font=\tiny, text=black!50, anchor=south, rotate=45]
    at (axis cs:2012,105) {AlexNet};

% Transformers
\draw[black!40, dashed, thin] (axis cs:2017,0) -- (axis cs:2017,105);
\node[font=\tiny, text=black!50, anchor=south, rotate=45]
    at (axis cs:2017,105) {Transformers};

% GPT-3
\draw[black!40, dashed, thin] (axis cs:2020,0) -- (axis cs:2020,105);
\node[font=\tiny, text=black!50, anchor=south, rotate=45]
    at (axis cs:2020,105) {GPT-3};

% GPT-4 / AutoGen
\draw[black!40, dashed, thin] (axis cs:2023.25,0) -- (axis cs:2023.25,105);
\node[font=\tiny, text=black!50, anchor=south, rotate=45]
    at (axis cs:2023.25,105) {GPT-4};

% ── Data series ────────────────────────────────────────────
% "Intelligent agent" — the dominant new term (area fill + line)
\addplot[
    color=blue!70!black,
    fill=blue!10,
    thick,
    smooth,
    mark=none,
] table[x index=0, y index=3] {figures/search-trends.dat}
    \closedcycle;

% "Agentic AI" — the new buzzword
\addplot[
    color=red!70!black,
    fill=red!10,
    thick,
    smooth,
    mark=none,
] table[x index=0, y index=2] {figures/search-trends.dat}
    \closedcycle;

% "Multi-agent system" — the classical term
\addplot[
    color=black!70!green,
    fill=green!8,
    thick,
    smooth,
    mark=none,
] table[x index=0, y index=1] {figures/search-trends.dat}
    \closedcycle;

\legend{
    ``Intelligent agent'',
    ``Agentic AI'',
    ``Multi-agent system''
}

% ── Peak annotation ────────────────────────────────────────
\node[font=\tiny, anchor=south west, text=blue!70!black]
    at (axis cs:2025.3,88) {100 (Oct 2025)};

% ── Classical era label ────────────────────────────────────
\draw[decorate, decoration={brace, amplitude=4pt, mirror},
      black!50, thick]
    (axis cs:2004,-2) -- (axis cs:2006,-2)
    node[midway, below=5pt, font=\tiny, text=black!50]
    {Classical MAS};

\end{axis}
\end{tikzpicture}
\caption{\textbf{The Attention Gap.} Google Trends search interest
(2004--2026) for three agent-related terms. ``Multi-agent system''
(the classical terminology) peaked at 3--4 in 2004 and flatlined near
zero by 2012. ``Intelligent agent'' and ``Agentic AI'' (new terminology)
exploded starting Q4 2024, reaching 100 by October 2025. The shaded
region marks the Attention Gap (2013--2022): a decade where agent research
effectively disappeared from public discourse. The new wave of interest
uses different terminology than the classical research it unknowingly
reinvents---the terminological disconnect that motivates this survey.
Data source: Google Trends, quarterly averages.}
\label{fig:trends}
\end{figure*}


\begin{figure*}[t]
\centering
\begin{tikzpicture}
\begin{axis}[
    width=\textwidth,
    height=7.5cm,
    xlabel={Year},
    ylabel={Google Trends Interest (0--100)},
    xmin=2004, xmax=2026.25,
    ymin=-3, ymax=108,
    xtick={2004,2006,2008,2010,2012,2014,2016,2018,2020,2022,2024,2026},
    xticklabel style={font=\small},
    yticklabel style={font=\small},
    xlabel style={font=\small},
    ylabel style={font=\small},
    legend style={
        at={(0.02,0.98)},
        anchor=north west,
        font=\small,
        draw=black!30,
        fill=white,
        fill opacity=0.9,
        text opacity=1,
        cells={anchor=west},
    },
    grid=major,
    grid style={gray!20},
    clip=false,
    % Era shading
    extra x ticks={},
]

% ── The Dead Zone shading (2013-2022) ──────────────────────
\fill[black!5] (axis cs:2012.5,-3) rectangle (axis cs:2022.5,108);
\node[font=\footnotesize\itshape, text=black!40, rotate=90]
    at (axis cs:2017.5,50) {The Attention Gap};

% ── Key event markers ─────────────────────────────────────
% AlexNet
\draw[black!40, dashed, thin] (axis cs:2012,0) -- (axis cs:2012,105);
\node[font=\tiny, text=black!50, anchor=south, rotate=45]
    at (axis cs:2012,105) {AlexNet};

% Transformers
\draw[black!40, dashed, thin] (axis cs:2017,0) -- (axis cs:2017,105);
\node[font=\tiny, text=black!50, anchor=south, rotate=45]
    at (axis cs:2017,105) {Transformers};

% GPT-3
\draw[black!40, dashed, thin] (axis cs:2020,0) -- (axis cs:2020,105);
\node[font=\tiny, text=black!50, anchor=south, rotate=45]
    at (axis cs:2020,105) {GPT-3};

% GPT-4 / AutoGen
\draw[black!40, dashed, thin] (axis cs:2023.25,0) -- (axis cs:2023.25,105);
\node[font=\tiny, text=black!50, anchor=south, rotate=45]
    at (axis cs:2023.25,105) {GPT-4};

% ── Data series ────────────────────────────────────────────
% "Intelligent agent" — the dominant new term (area fill + line)
\addplot[
    color=blue!70!black,
    fill=blue!10,
    thick,
    smooth,
    mark=none,
] table[x index=0, y index=3] {figures/search-trends.dat}
    \closedcycle;

% "Agentic AI" — the new buzzword
\addplot[
    color=red!70!black,
    fill=red!10,
    thick,
    smooth,
    mark=none,
] table[x index=0, y index=2] {figures/search-trends.dat}
    \closedcycle;

% "Multi-agent system" — the classical term
\addplot[
    color=black!70!green,
    fill=green!8,
    thick,
    smooth,
    mark=none,
] table[x index=0, y index=1] {figures/search-trends.dat}
    \closedcycle;

\legend{
    ``Intelligent agent'',
    ``Agentic AI'',
    ``Multi-agent system''
}

% ── Peak annotation ────────────────────────────────────────
\node[font=\tiny, anchor=south west, text=blue!70!black]
    at (axis cs:2025.3,88) {100 (Oct 2025)};

% ── Classical era label ────────────────────────────────────
\draw[decorate, decoration={brace, amplitude=4pt, mirror},
      black!50, thick]
    (axis cs:2004,-2) -- (axis cs:2006,-2)
    node[midway, below=5pt, font=\tiny, text=black!50]
    {Classical MAS};

\end{axis}
\end{tikzpicture}
\caption{\textbf{The Attention Gap.} Google Trends search interest
(2004--2026) for three agent-related terms. ``Multi-agent system''
(the classical terminology) peaked at 3--4 in 2004 and flatlined near
zero by 2012. ``Intelligent agent'' and ``Agentic AI'' (new terminology)
exploded starting Q4 2024, reaching 100 by October 2025. The shaded
region marks the Attention Gap (2013--2022): a decade where agent research
effectively disappeared from public discourse. The new wave of interest
uses different terminology than the classical research it unknowingly
reinvents---the terminological disconnect that motivates this survey.
Data source: Google Trends, quarterly averages.}
\label{fig:trends}
\end{figure*}


\begin{figure*}[t]
\centering
\begin{tikzpicture}
\begin{axis}[
    width=\textwidth,
    height=7.5cm,
    xlabel={Year},
    ylabel={Google Trends Interest (0--100)},
    xmin=2004, xmax=2026.25,
    ymin=-3, ymax=108,
    xtick={2004,2006,2008,2010,2012,2014,2016,2018,2020,2022,2024,2026},
    xticklabel style={font=\small},
    yticklabel style={font=\small},
    xlabel style={font=\small},
    ylabel style={font=\small},
    legend style={
        at={(0.02,0.98)},
        anchor=north west,
        font=\small,
        draw=black!30,
        fill=white,
        fill opacity=0.9,
        text opacity=1,
        cells={anchor=west},
    },
    grid=major,
    grid style={gray!20},
    clip=false,
    % Era shading
    extra x ticks={},
]

% ── The Dead Zone shading (2013-2022) ──────────────────────
\fill[black!5] (axis cs:2012.5,-3) rectangle (axis cs:2022.5,108);
\node[font=\footnotesize\itshape, text=black!40, rotate=90]
    at (axis cs:2017.5,50) {The Attention Gap};

% ── Key event markers ─────────────────────────────────────
% AlexNet
\draw[black!40, dashed, thin] (axis cs:2012,0) -- (axis cs:2012,105);
\node[font=\tiny, text=black!50, anchor=south, rotate=45]
    at (axis cs:2012,105) {AlexNet};

% Transformers
\draw[black!40, dashed, thin] (axis cs:2017,0) -- (axis cs:2017,105);
\node[font=\tiny, text=black!50, anchor=south, rotate=45]
    at (axis cs:2017,105) {Transformers};

% GPT-3
\draw[black!40, dashed, thin] (axis cs:2020,0) -- (axis cs:2020,105);
\node[font=\tiny, text=black!50, anchor=south, rotate=45]
    at (axis cs:2020,105) {GPT-3};

% GPT-4 / AutoGen
\draw[black!40, dashed, thin] (axis cs:2023.25,0) -- (axis cs:2023.25,105);
\node[font=\tiny, text=black!50, anchor=south, rotate=45]
    at (axis cs:2023.25,105) {GPT-4};

% ── Data series ────────────────────────────────────────────
% "Intelligent agent" — the dominant new term (area fill + line)
\addplot[
    color=blue!70!black,
    fill=blue!10,
    thick,
    smooth,
    mark=none,
] table[x index=0, y index=3] {figures/search-trends.dat}
    \closedcycle;

% "Agentic AI" — the new buzzword
\addplot[
    color=red!70!black,
    fill=red!10,
    thick,
    smooth,
    mark=none,
] table[x index=0, y index=2] {figures/search-trends.dat}
    \closedcycle;

% "Multi-agent system" — the classical term
\addplot[
    color=black!70!green,
    fill=green!8,
    thick,
    smooth,
    mark=none,
] table[x index=0, y index=1] {figures/search-trends.dat}
    \closedcycle;

\legend{
    ``Intelligent agent'',
    ``Agentic AI'',
    ``Multi-agent system''
}

% ── Peak annotation ────────────────────────────────────────
\node[font=\tiny, anchor=south west, text=blue!70!black]
    at (axis cs:2025.3,88) {100 (Oct 2025)};

% ── Classical era label ────────────────────────────────────
\draw[decorate, decoration={brace, amplitude=4pt, mirror},
      black!50, thick]
    (axis cs:2004,-2) -- (axis cs:2006,-2)
    node[midway, below=5pt, font=\tiny, text=black!50]
    {Classical MAS};

\end{axis}
\end{tikzpicture}
\caption{\textbf{The Attention Gap.} Google Trends search interest
(2004--2026) for three agent-related terms. ``Multi-agent system''
(the classical terminology) peaked at 3--4 in 2004 and flatlined near
zero by 2012. ``Intelligent agent'' and ``Agentic AI'' (new terminology)
exploded starting Q4 2024, reaching 100 by October 2025. The shaded
region marks the Attention Gap (2013--2022): a decade where agent research
effectively disappeared from public discourse. The new wave of interest
uses different terminology than the classical research it unknowingly
reinvents---the terminological disconnect that motivates this survey.
Data source: Google Trends, quarterly averages.}
\label{fig:trends}
\end{figure*}
